\section*{Hoofdstuk \ref{ch:g4:world-physical-map} --- De natuurkundige kaart van de wereld}

\begin{solution}{exo:g4:world-physical-map:1}
De \omterm{def:g4:world-physical-map:altitude}{hoogte} van een plek is hoe hoog die boven \omterm{def:g4:world-physical-map:altitude}{zeeniveau} ligt, in meter. Ze
wordt gemeten vanaf het \omterm{def:g4:world-physical-map:altitude}{zeeniveau}, het gemiddelde niveau van het
zeeoppervlak.
\end{solution}

\begin{solution}{exo:g4:world-physical-map:2}
(a) Donkergroen: tussen 0 en 200 meter. (b) Lichtgeel: tussen 500 en
1\,000 meter. (c) Bijna wit: hoger dan 4\,000 meter.
\end{solution}

\begin{solution}{exo:g4:world-physical-map:3}
(a) Zuid-Amerika. (b) Azië. (c) Afrika. (d) Noord-Amerika.
\end{solution}

\begin{solution}{exo:g4:world-physical-map:4}
De hoogste \omterm{def:g1:landscapes-around-us:landform}{berg} is Mount Everest (8\,849 meter). Het grootste meer is de
Kaspische Zee; het diepste meer is het Baikalmeer.
\end{solution}

\begin{solution}{exo:g4:world-physical-map:5}
Hoe donkerder het blauw, hoe dieper de \omterm{def:g2:water-on-earth:ocean}{zee}: bij de \omterm{def:g1:landscapes-around-us:water}{kusten} is de \omterm{def:g2:water-on-earth:ocean}{zee} ondiep
(minder dan 200 meter) en midden in de \omterm{def:g2:water-on-earth:ocean}{oceanen} veel dieper (op veel plaatsen
meer dan 4\,000 meter).
\end{solution}

\begin{solution}{exo:g4:world-physical-map:6}
De \omterm{def:g4:world-physical-map:physical-map}{natuurkundige kaart} laat de \omterm{def:g4:world-physical-map:altitude}{hoogte} van het oppervlak zien. Op Antarctica
en Groenland is het oppervlak de bovenkant van een dikke laag ijs, die hoog
boven \omterm{def:g4:world-physical-map:altitude}{zeeniveau} ligt; daarom krijgen ze de kleuren van hoog land.
\end{solution}

\begin{solution}{exo:g4:world-physical-map:7}
De bodem van de Challengerdiepte (bijna 11\,000 meter onder \omterm{def:g4:world-physical-map:altitude}{zeeniveau}); de
gemiddelde oceaanbodem (ongeveer 3\,700 meter eronder); een strand
(ongeveer 0 meter); een dorp op 1\,500 meter; de top van Mount Everest
(8\,849 meter).
\end{solution}

\begin{solution}{exo:g4:world-physical-map:8}
Een \omterm{def:g1:landscapes-around-us:water}{rivier} bestaat uit veel beken die samenkomen, en men is het er niet
over eens welke beek haar echte begin is. Gemeten vanaf een verder gelegen
bron is een \omterm{def:g1:landscapes-around-us:water}{rivier} langer. De Nijl en de Amazone zijn bijna even lang (elk
ongeveer 6\,700 tot 6\,800 kilometer), dus het antwoord hangt af van die
keuze.
\end{solution}

\begin{solution}{exo:g4:world-physical-map:9}
Een kwart van 30 miljoen is zevenenhalf miljoen vierkante kilometer (30
gedeeld door 4); een derde is 10 miljoen (30 gedeeld door 3). De Sahara,
achtenhalf miljoen, is meer dan een kwart van Afrika maar minder dan een
derde.
\end{solution}

\begin{solution}{exo:g4:world-physical-map:10}
$2500 - 1000 = 1500$: de hut ligt 1\,500 meter hoger. De hut, tussen
2\,000 en 4\,000 meter, is bruin; het \omterm{def:g1:landscapes-around-us:water}{meer}, op 1\,000 meter, ligt op de grens
tussen lichtgeel (500 tot 1\,000) en lichtbruin (1\,000 tot 2\,000): beide
worden goedgerekend.
\end{solution}

\begin{solution}{exo:g4:world-physical-map:11}
10 kilometer is 10\,000 meter. $10000 - 8849 = 1151$: het vliegtuig vliegt
hoger dan Mount Everest, 1\,151 meter hoger. Boven de Challengerdiepte
vliegt het 10\,000 meter boven de \omterm{def:g2:water-on-earth:ocean}{zee}, en de bodem van de trog ligt
11\,000 meter onder de \omterm{def:g2:water-on-earth:ocean}{zee}: $10000 + 11000 = 21000$ meter boven de bodem
van de trog.
\end{solution}
